\section{Introduction}

A conversation can end while the work continues. A maintainer leaves, an
agent process stops, a repository moves, a release is promoted, or a new
participant arrives with no access to the private context that made the last
decision seem obvious. If the work is to survive, the newcomer must recover
more than a summary. The newcomer must be able to identify what happened,
decide which claims may be relied upon, and determine whether and how to
continue.

This is a systems problem before it is an interface problem. Better models,
larger context windows, and more persuasive explanations do not by themselves
preserve the identity of the work across time. Nor does a complete event log
automatically tell a participant what a claim proves or why cooperation is
worth entering. Persistent work needs a shared world that can outlive the
participant who currently interprets it.

The first three Kung Fu Decisions (KFDs) state a compact foundation for that
world \cite{kfd2026alpha42}:

\begin{quote}
KFD-1: facts must not drift.\\
KFD-2: trust must start from facts.\\
KFD-3: cooperation must start from trusted value.
\end{quote}

Read as three imperatives, the statements can sound like good advice. This
paper advances a stronger and more falsifiable interpretation: they name three
different roles that persistent cooperation cannot collapse into one another.
KFD-1 preserves the continuity of the world about which participants reason.
KFD-2 converts parts of that world into bounded permission to rely. KFD-3
converts warranted reliance and visible value into coordinated action among
participants who retain independent judgment.

The dependency is ordered, but its operation is recursive:

\begin{equation}
\text{facts} \rightarrow \text{trust} \rightarrow \text{cooperation}
\rightarrow \text{new occurrences} \rightarrow \text{successor facts}.
\label{eq:intro-loop}
\end{equation}

This loop is the paper's central object. It explains why the triad can be
foundational without being a complete ontology. The three KFDs do not name
every object a cooperative system will need. They preserve the conditions
under which new action can occur, its consequences can correct the shared
world, and missing objects can later be discovered without replacing reality
with narrative.

The paper makes four contributions:

\begin{itemize}[leftmargin=*]
  \item it defines facts, trust, and cooperation as operational roles rather
        than moral qualities;
  \item it uses removal and order-inversion tests to explain why the roles are
        separately necessary and directionally dependent;
  \item it models the triad as a recursive kernel through which a shared work
        world persists, becomes actionable, and produces further evidence;
        and
  \item it derives a narrow frontier from that kernel: declared perspective
        in KFD-4, separated genesis and qualification in KFD-5, and the
        causal-experience gate in the draft KFD-6.
\end{itemize}

The claims remain bounded. The paper does not prove that the triad is a
mathematically unique basis for every social or technical system. It does not
equate recorded claims with objective reality, trust with certainty, or
cooperation with agreement. It argues that any system claiming durable
cooperation must implement equivalent roles or explain how it preserves the
same continuity, warranted reliance, and generative coordination without
them.

Section~\ref{sec:problem} defines the continuity problem.
Section~\ref{sec:foundation} specifies the three roles.
Section~\ref{sec:necessity} applies removal and order tests.
Section~\ref{sec:kernel} develops the recursive model and a cross-session
case. Section~\ref{sec:frontier} shows how KFD-4/5/6 arise at the first
pressure frontier. Sections~\ref{sec:related} and~\ref{sec:discussion}
position and bound the claim.
